
\sgaStableTarget{sga6.e09.prop.2-2}\textbf{Proposition 2.2.} \uline{With the notation and hypotheses of} 2.1, \uline{let} @@M0@@ \uline{be the structural morphism of} @@M1@@. \uline{Then the homomorphism}:
@@M2@@
\uline{is surjective}, \uline{and its kernel is the subgroup of} @@M3@@ \uline{generated by the subgroups} @@M4@@ \uline{for} @@M5@@, \uline{or}, \uline{equivalently}, \uline{for} @@M6@@ \uline{running through a system of generators of} @@M7@@.

First suppose @@M8@@, i.e.\ @@M9@@. Then @@M10@@, and for every @@M11@@ one has:
@@M12@@
Consequently, putting @@M13@@, one has @@M14@@, which shows that @@M15@@ identifies with the complement of the zero section in the bundle @@M16@@. Considering @@M17@@ as a closed subscheme of @@M18@@ by the zero section, one has, by 1.1, the exact sequence:
@@M19@@
where @@M20@@ is the zero section of @@M21@@. Moreover, @@M22@@ factors as @@M23@@, where @@M24@@ is the structural morphism of @@M25@@. By 1.6, @@M26@@ is an isomorphism, inverse to @@M27@@. Thus one has the exact sequence:
@@M28@@
from which one obtains, by 1.8:
@@M29@@
where @@M30@@ is the class in @@M31@@ of @@M32@@. Since for every @@M33@@, @@M34@@ is a multiple of @@M35@@, this proves 2.2 when @@M36@@ has rank @@M37@@.

